\documentclass[11pt,a4paper]{article}

\usepackage[margin=1in]{geometry}
\usepackage[T1]{fontenc}
\usepackage[utf8]{inputenc}
\usepackage{lmodern}
\usepackage{hyperref}
\usepackage{enumitem}
\usepackage{booktabs}
\usepackage{longtable}
\usepackage{xcolor}
\usepackage{listings}

\hypersetup{
  colorlinks=true,
  linkcolor=blue!60!black,
  urlcolor=blue!60!black
}

\lstset{
  basicstyle=\ttfamily\small,
  breaklines=true,
  frame=single,
  columns=fullflexible
}

\title{ERC Indoor Progress Note: March 19\\
\large What We Built, Why We Built It, and Where We Are Now}
\author{Working Project Note}
\date{\today}

\begin{document}

\maketitle

\section{What This Document Is}

This document explains, in a simple but technically honest way, what has been
built so far for the ERC indoor project inside this repository.  The goal is
not just to list files.  The goal is to explain:

\begin{itemize}[nosep]
  \item what each part does,
  \item why we chose it,
  \item what alternatives existed,
  \item why those alternatives were not chosen as the main baseline,
  \item and where the project stands right now.
\end{itemize}

The intended result is that after reading this, a teammate should understand the
current project structure clearly enough to contribute to it and defend the main
engineering choices.


\section{The Actual Problem We Are Solving}

We are building an \textbf{indoor navigation system} for the
\textbf{EarthRover Challenge} in a \textbf{known and repeated corridor
environment}.

That last point matters a lot.

This is \textbf{not} the same as general indoor navigation in an unknown
building.  We are not trying to make a robot that can enter any building in the
world and navigate from scratch.  We are solving a more specific and more
practical problem:

\begin{itemize}[nosep]
  \item the corridor or building family is known,
  \item we can collect data in advance,
  \item we can test in the same environment repeatedly,
  \item and the competition checkpoints are image-based, not GPS-based.
\end{itemize}

This changes the whole strategy.  It means we are allowed to build a
corridor-specific baseline and tune it carefully for that environment.


\section{Why The Final Direction Changed}

At the start, many possible directions existed:

\begin{itemize}[nosep]
  \item metric SLAM,
  \item pure visual odometry,
  \item a large end-to-end controller,
  \item CityWalker-style reuse,
  \item MBRA as the whole stack,
  \item or a graph-based visual localization system.
\end{itemize}

After looking at the repo and the actual competition setting, the most sensible
backbone became:

\begin{center}
\fbox{\parbox{0.9\textwidth}{
\centering
\textbf{known-corridor visual localization + graph planning + local controller}
}}
\end{center}

Why this direction won:

\begin{itemize}[nosep]
  \item goals are image-defined,
  \item the corridor is repeated and known,
  \item monocular RGB is available,
  \item metric localization is weak,
  \item latency and frame rate are poor,
  \item and we needed something debuggable and reliable.
\end{itemize}

This does \emph{not} mean the other ideas are useless.  It means they are not
the cleanest \textbf{main backbone} for this exact problem.


\section{The Most Important Decision}

The most important decision we made is:

\begin{quote}
Use a \textbf{corridor-specific perception/planning backbone} first, then plug
in a controller on top of it.
\end{quote}

That means:

\begin{itemize}[nosep]
  \item first answer ``where am I?'',
  \item then answer ``where should I go next?'',
  \item then answer ``how do I move there safely?''.
\end{itemize}

We did \textbf{not} start from the controller and hope it solves everything.


\section{What We Built}

\subsection{1. A Clean Data Pipeline From The Teleoperation H5 File}

We found that the teleoperation recording file
\texttt{data/corrider.h5} is extremely useful.

It contains:

\begin{itemize}[nosep]
  \item front camera frames,
  \item control commands,
  \item telemetry,
  \item orientation,
  \item gyros,
  \item accelerometers,
  \item magnetometer data,
  \item and RPM samples.
\end{itemize}

This was better than expected.  It means we do not need to start by collecting
everything again from zero.

To make this usable, we built:

\begin{itemize}[nosep]
  \item \texttt{tools/extract\_h5\_dataset.py}
\end{itemize}

What it does:

\begin{itemize}[nosep]
  \item extracts front frames into an ordered image folder,
  \item exports CSV metadata,
  \item builds a baseline-friendly \texttt{data\_info.json},
  \item and aligns the different streams using relative time.
\end{itemize}

Why this helps:

\begin{itemize}[nosep]
  \item the localization system wants images,
  \item the planning system wants ordered steps,
  \item and the control / state-estimation side needs actions and motion
        signals.
\end{itemize}

Why we chose this over ad hoc notebooks:

\begin{itemize}[nosep]
  \item it is reusable,
  \item it is shareable,
  \item and it reduces the chance of hidden one-off mistakes.
\end{itemize}


\subsection{2. A Real Baseline Database Builder}

We turned the earlier notebook-style corridor work into a usable script:

\begin{itemize}[nosep]
  \item \texttt{baseline.py}
\end{itemize}

What it does:

\begin{itemize}[nosep]
  \item extracts CosPlace descriptors,
  \item builds a place graph,
  \item builds a navigation/action graph when \texttt{data\_info.json} is
        available,
  \item and supports query-time localization.
\end{itemize}

We used it to build:

\begin{itemize}[nosep]
  \item \texttt{data/corrider\_db/}
\end{itemize}

This contains:

\begin{itemize}[nosep]
  \item \texttt{descriptors.npz}
  \item \texttt{config.json}
  \item \texttt{place\_graph.json}
  \item \texttt{navigation\_graph.json}
\end{itemize}

Why this matters:

\begin{itemize}[nosep]
  \item this is the corridor memory of the system,
  \item without it, we do not have a known-environment localization backbone.
\end{itemize}

Why we chose CosPlace + graph:

\begin{itemize}[nosep]
  \item it matches the known-corridor setting,
  \item it is explainable,
  \item it can be debugged visually,
  \item and it is already close to the older project code direction.
\end{itemize}

Alternative option:

\begin{itemize}[nosep]
  \item use ORB-SLAM3 or direct SLAM as the main backbone
\end{itemize}

Why we did not choose that as the baseline:

\begin{itemize}[nosep]
  \item the problem is image-goal and corridor-specific, not metric-map-first,
  \item repeated corridor appearance is risky for monocular SLAM,
  \item and the graph baseline better matches what we already had.
\end{itemize}


\subsection{3. Temporal Localization}

Single-frame retrieval is not enough in a repeated corridor.  Even when the
database is good, one frame can still be ambiguous.

So we built:

\begin{itemize}[nosep]
  \item \texttt{src/temporal\_localization.py}
\end{itemize}

What it does:

\begin{itemize}[nosep]
  \item takes top retrieval candidates,
  \item penalizes impossible jumps,
  \item penalizes unreasonable backwards motion,
  \item optionally uses heading consistency,
  \item and resolves ambiguity conservatively.
\end{itemize}

Why it helps:

\begin{itemize}[nosep]
  \item a robot does not teleport,
  \item so localization should respect continuity,
  \item and repeated corridor segments need temporal context.
\end{itemize}

We then built:

\begin{itemize}[nosep]
  \item \texttt{tools/evaluate\_temporal\_localization.py}
\end{itemize}

This let us test temporal localization on held-out frames.

The most important outcome:

\begin{itemize}[nosep]
  \item overall exact match was strong,
  \item but more importantly, moving-frame performance was extremely strong,
  \item and the worst-looking error cluster turned out to be mostly a repeated
        stationary segment, not true moving localization failure.
\end{itemize}

This is a big result because it means the perception/planning backbone is
actually stronger than the naive error count first suggested.


\subsection{4. Runtime Corridor Localizer}

The database and temporal localizer were not enough by themselves because they
were still pieces, not a runtime API.

So we built:

\begin{itemize}[nosep]
  \item \texttt{src/corridor\_localizer.py}
\end{itemize}

What it does:

\begin{itemize}[nosep]
  \item loads the built DB once,
  \item encodes a new frame,
  \item retrieves candidates,
  \item applies temporal localization,
  \item and returns a planning-ready result:
        node, step, confidence, and candidate list.
\end{itemize}

Why this helps:

\begin{itemize}[nosep]
  \item planning should not need to know how descriptor extraction works,
  \item it should just ask: ``where are we right now?''.
\end{itemize}


\subsection{5. Runtime Graph Planner}

Once localization exists, the next question is:

\begin{quote}
Given the current node and a target checkpoint, what nearby subgoal should the
controller follow right now?
\end{quote}

So we built:

\begin{itemize}[nosep]
  \item \texttt{src/graph\_planner.py}
\end{itemize}

What it does:

\begin{itemize}[nosep]
  \item loads the runtime graph,
  \item resolves targets by node, step, or image name,
  \item computes shortest path,
  \item chooses a nearby subgoal node,
  \item and supports checkpoint progression.
\end{itemize}

Why it helps:

\begin{itemize}[nosep]
  \item the controller should not chase the final goal from far away,
  \item it should chase a nearby subgoal on a valid corridor path.
\end{itemize}

Alternative option:

\begin{itemize}[nosep]
  \item let a learned controller implicitly handle the whole route
\end{itemize}

Why we did not choose that as the baseline:

\begin{itemize}[nosep]
  \item it would be much harder to debug,
  \item and the graph planner already gives us route structure cleanly.
\end{itemize}


\subsection{6. Controller-Facing Runtime Coordinator}

At that point, we had localization and planning, but control still had to
manually call several pieces.

So we built:

\begin{itemize}[nosep]
  \item \texttt{src/navigation\_runtime.py}
\end{itemize}

What it does:

\begin{itemize}[nosep]
  \item combines localization and planning,
  \item returns a controller-facing bundle,
  \item and can optionally load the subgoal image for the controller.
\end{itemize}

Why this helps:

\begin{itemize}[nosep]
  \item control should receive one clean package:
        current state, target, subgoal, confidence, and subgoal image.
\end{itemize}


\subsection{7. A First Local Controller Baseline}

We still needed a controller that can turn a nearby subgoal into a motion
command.

So we built:

\begin{itemize}[nosep]
  \item \texttt{src/local\_controller.py}
\end{itemize}

What it does:

\begin{itemize}[nosep]
  \item takes the controller bundle,
  \item uses current step, subgoal step, current orientation, subgoal
        orientation, and confidence,
  \item and outputs a simple baseline:
        \texttt{linear}, \texttt{angular}, and a reason.
\end{itemize}

Why this helps:

\begin{itemize}[nosep]
  \item it gives us a real executable baseline before MBRA integration,
  \item it is simple enough to debug,
  \item and it gives the control side a starting point immediately.
\end{itemize}

Why this is not the final answer:

\begin{itemize}[nosep]
  \item it is heuristic,
  \item not learned,
  \item and not yet safety-aware.
\end{itemize}

But it is still valuable because it proves the runtime pipeline can already
produce commands.


\subsection{8. A Conservative Live Runner}

Finally, we connected the pieces into a live loop:

\begin{itemize}[nosep]
  \item \texttt{live\_indoor\_runtime.py}
\end{itemize}

What it does:

\begin{itemize}[nosep]
  \item connects to the SDK,
  \item gets camera and heading,
  \item runs localization,
  \item runs planning,
  \item runs the local controller,
  \item and optionally sends control to the robot.
\end{itemize}

Important safety choice:

\begin{itemize}[nosep]
  \item it defaults to \textbf{dry-run},
  \item and only sends real control if explicitly requested.
\end{itemize}

Why this matters:

\begin{itemize}[nosep]
  \item we should not accidentally move the robot during first tests,
  \item especially when the runtime stack is still being validated.
\end{itemize}


\subsection{9. A Lightweight Motion-Prior Layer}

One thing that was missing in the first version of the stack was a proper
middle layer between raw robot telemetry and the controller.

So we added:

\begin{itemize}[nosep]
  \item \texttt{src/sensor\_state.py}
\end{itemize}

What it does:

\begin{itemize}[nosep]
  \item smooths heading from SDK orientation,
  \item estimates turn rate from gyro $z$,
  \item summarizes RPMs into a weak motion hint,
  \item and passes this filtered motion state into localization and control.
\end{itemize}

Why this matters:

\begin{itemize}[nosep]
  \item a live robot should not use raw noisy heading directly if we can avoid
        it,
  \item gyro information helps tell whether the robot is already turning,
  \item and this makes the controller less twitchy.
\end{itemize}

Why we chose this instead of jumping straight to an EKF:

\begin{itemize}[nosep]
  \item the visual place-localization backbone is already the main state
        estimate,
  \item we do not yet have a trustworthy metric motion model for the whole
        indoor problem,
  \item and a small filter gives us most of the immediate benefit without
        pretending we solved full state estimation.
\end{itemize}

So the correct interpretation is:

\begin{quote}
We are using a \textbf{lightweight motion prior}, not replacing the visual
backbone with a Kalman-filtered metric localization system.
\end{quote}


\section{What About IMU, Orientation, and RPMs?}

One correction that became important during this work:

\begin{quote}
The project should not be described as vision-only.
\end{quote}

The correct statement is:

\begin{quote}
\textbf{vision-primary with IMU / heading support}
\end{quote}

What that means:

\begin{itemize}[nosep]
  \item camera frames are still the main localization input,
  \item but orientation and IMU should help temporal filtering, recovery, and
        controller smoothing,
  \item and RPM data should be used only if it proves trustworthy on the real
        robot.
\end{itemize}

Why we did not make IMU the main baseline:

\begin{itemize}[nosep]
  \item IMU alone drifts,
  \item the task is image-goal and corridor-specific,
  \item and the strongest existing code already centered on visual place-based
        reasoning.
\end{itemize}

So the right role of IMU is:

\begin{itemize}[nosep]
  \item support signal,
  \item not sole localization backbone.
\end{itemize}

The practical version of that sentence is:

\begin{itemize}[nosep]
  \item filtered heading helps the localizer,
  \item gyro turn rate helps the controller,
  \item RPMs can give weak motion hints,
  \item but visual localization still decides where we are in the corridor.
\end{itemize}


\section{What About MBRA?}

This is one of the most important conceptual clarifications.

The short answer is:

\begin{quote}
\texttt{mbra\_repo} is \textbf{not} the current deployed indoor backbone.
\end{quote}

What it is:

\begin{itemize}[nosep]
  \item a research reference repo,
  \item containing MBRA / LogoNav-related model and deployment code,
  \item useful as a candidate controller direction.
\end{itemize}

What it is \textbf{not} currently:

\begin{itemize}[nosep]
  \item not the active perception/planning backbone,
  \item not fully integrated into our corridor runtime,
  \item not already a clean indoor drop-in controller in this project.
\end{itemize}

Why we did not force MBRA into the stack immediately:

\begin{itemize}[nosep]
  \item the corridor perception/planning side needed to be solid first,
  \item and the repo itself does not automatically provide a finished indoor
        controller plug-in for our exact setup.
\end{itemize}

So the current status of MBRA is:

\begin{itemize}[nosep]
  \item important candidate for the controller side,
  \item not yet the active controller in the baseline loop.
\end{itemize}


\section{Localization Explained Very Clearly}

This is the most important part to understand, because this is the part that is
currently working well.

\subsection{What localization means here}

Localization means:

\begin{quote}
Given the current live camera frame, estimate which stored corridor frame or
graph node the robot is closest to.
\end{quote}

We are \textbf{not} doing metric $(x, y)$ localization on a floorplan.
We are doing \textbf{place localization in a known visual corridor graph}.

\subsection{What information goes in}

The localization pipeline takes:

\begin{itemize}[nosep]
  \item a live RGB frame from the robot,
  \item the built database of reference corridor images,
  \item the descriptor vectors for those reference images,
  \item optional heading information,
  \item and temporal history from previous frames.
\end{itemize}

\subsection{What information comes out}

The localization pipeline outputs:

\begin{itemize}[nosep]
  \item the best-matching graph node,
  \item the corresponding corridor step,
  \item a confidence score,
  \item and a shortlist of alternative candidates.
\end{itemize}

\subsection{The actual logic}

The logic has three layers.

\paragraph{Layer 1: image embedding.}

Each reference image is passed through CosPlace and converted into a feature
vector of dimension $512$.

Call the reference descriptors:
\[
d_1, d_2, \dots, d_N
\]

and the live query descriptor:
\[
q
\]

The first retrieval step asks:

\begin{quote}
Which stored descriptor $d_i$ is most similar to $q$?
\end{quote}

That gives top candidate nodes.

\paragraph{Layer 2: retrieval score.}

At the simplest level, retrieval is just nearest-neighbor search in descriptor
space.

Conceptually, we rank candidates by a similarity score such as:
\[
\mathrm{sim}(q, d_i)
\]

Higher similarity means the current live view looks more like reference frame
$i$.

\paragraph{Layer 3: temporal filtering.}

This is the part that made the system actually usable.

A single frame can be ambiguous in a corridor.  So we do not trust only the raw
best match.  We also penalize unreasonable jumps.

Conceptually, the score becomes:
\[
\mathrm{score}(i) =
\mathrm{sim}(q, d_i)
- \lambda_{\text{jump}} \cdot \mathrm{jump\_cost}
- \lambda_{\text{back}} \cdot \mathrm{backward\_cost}
- \lambda_{\text{heading}} \cdot \mathrm{heading\_cost}
\]

This is not meant as theory for theory's sake.  It is exactly the practical
idea implemented in \texttt{src/temporal\_localization.py}.

In plain terms:

\begin{itemize}[nosep]
  \item if a candidate would require the robot to teleport many steps, punish it,
  \item if a candidate implies weird backward motion, punish it,
  \item if a candidate disagrees strongly with heading, punish it,
  \item if two candidates are too close in score, prefer stability over jumping.
\end{itemize}

\subsection{Why this worked well}

Localization worked well for us because several things lined up correctly:

\begin{itemize}[nosep]
  \item the corridor is known and repeated,
  \item the database was built from the same environment family,
  \item CosPlace is good at place-level visual retrieval,
  \item temporal filtering prevents silly jumps,
  \item and the graph/node representation keeps the problem simple.
\end{itemize}

This is why the system can look at the live frame and say:

\begin{quote}
``This is probably step 313, not step 850 and not step 40.''
\end{quote}

\subsection{What exactly we proved}

We already proved the following:

\begin{itemize}[nosep]
  \item the teleop H5 can be converted into a clean localization dataset,
  \item the corridor database can be built,
  \item the localizer retrieves the correct corridor neighborhood,
  \item temporal filtering makes it stable,
  \item and in live dry-run testing, localization can lock to the correct region
        when the robot is placed in the corridor.
\end{itemize}

So the current strongest statement is:

\begin{quote}
The \textbf{place-localization backbone is real and useful}.  This is no longer
the vague part of the project.
\end{quote}


\section{Planning Explained Very Clearly}

Now that localization is working, the next question is:

\begin{quote}
Once we know where we are, how do we decide where to go next?
\end{quote}

\subsection{What planning means here}

Planning here does \textbf{not} mean occupancy-grid A* on a metric map.

Planning here means:

\begin{quote}
Take the current localized graph node and the desired target node, then compute
a path through the corridor graph and choose a nearby subgoal.
\end{quote}

\subsection{What goes into planning}

The planner needs:

\begin{itemize}[nosep]
  \item current node from localization,
  \item current step,
  \item target step or target node,
  \item graph connectivity from \texttt{navigation\_graph.json}.
\end{itemize}

\subsection{What comes out of planning}

The planner outputs:

\begin{itemize}[nosep]
  \item a path from current node to target node,
  \item a path in step numbers,
  \item a subgoal node a few hops ahead,
  \item and the corresponding subgoal image.
\end{itemize}

\subsection{Why we need a subgoal}

We do not want the controller to chase frame $400$ directly from frame $313$ as
one giant jump.  That is too much.

Instead, if the path is:
\[
313 \rightarrow 314 \rightarrow 315 \rightarrow 316 \rightarrow \dots \rightarrow 400
\]

the planner picks a nearby subgoal such as step $316$.

So planning is doing:

\begin{quote}
``Do not think about the whole route at once.  Think about the next small piece
of the route.''
\end{quote}

\subsection{What is right and wrong in planning right now}

What is already right:

\begin{itemize}[nosep]
  \item graph construction works,
  \item shortest-path planning works,
  \item subgoal selection works,
  \item and checkpoint-style planning structure exists.
\end{itemize}

What is still weak:

\begin{itemize}[nosep]
  \item the graph is mainly one-way because it follows teleop order,
  \item planner recovery when there is no valid path is still basic,
  \item and control handoff is stronger than the controller that receives it.
\end{itemize}

\subsection{What planning needs next}

The next planning-side work is not inventing a new planner from scratch.

It is:

\begin{itemize}[nosep]
  \item making target selection cleaner,
  \item handling no-path cases more gracefully,
  \item deciding whether we want bidirectional corridor edges,
  \item and integrating checkpoint progression more carefully.
\end{itemize}


\section{What Needs To Happen Next}

Now that localization is the strongest part, the next weak part is the
\textbf{controller and runtime behavior}.

The order of work should be:

\begin{enumerate}[nosep]
  \item verify that manual teleoperation works cleanly,
  \item improve the simple controller so it stops spinning in place,
  \item keep localization and planning fixed while debugging control,
  \item add safer runtime behavior for low confidence and no-path cases,
  \item only after that, decide whether MBRA should replace the simple controller.
\end{enumerate}

So the main project is no longer:

\begin{quote}
``Can we localize at all?''
\end{quote}

The main project is now:

\begin{quote}
``Can we turn good localization and good subgoals into stable robot motion?''
\end{quote}


\section{Why We Chose This Path Instead of Other Paths}

\subsection{Why not pure SLAM as the main backbone?}

\begin{itemize}[nosep]
  \item The task is image-goal based.
  \item The corridor is known, so corridor-specific place recognition is very
        attractive.
  \item The existing graph localization work was already closer to the actual
        problem.
\end{itemize}

\subsection{Why not a giant end-to-end policy first?}

\begin{itemize}[nosep]
  \item It would be harder to debug.
  \item We already had enough structure to do better than that.
  \item Reliability matters more than elegance here.
\end{itemize}

\subsection{Why not MBRA as the entire stack from day one?}

\begin{itemize}[nosep]
  \item MBRA is not the same thing as the full runtime system.
  \item The repo evidence suggested MBRA/LogoNav needed careful interpretation,
        not blind reuse.
  \item The graph-localization side was already much more concrete for this
        corridor task.
\end{itemize}

\subsection{Why this current path makes sense}

\begin{itemize}[nosep]
  \item It uses the known corridor advantage properly.
  \item It is modular.
  \item It is debuggable.
  \item It gives each teammate a clear layer to work on.
  \item It lets us plug in a better controller later without rebuilding the
        whole stack.
\end{itemize}


\section{Where We Are Right Now}

Right now, the project has:

\begin{itemize}[nosep]
  \item a usable dataset extraction pipeline,
  \item a built corridor database,
  \item a working temporal localizer,
  \item a working runtime localizer,
  \item a working runtime graph planner,
  \item a controller-facing coordinator,
  \item a first local-controller baseline,
  \item and a live loop script that can run in dry-run mode.
\end{itemize}

In one line:

\begin{center}
\fbox{\parbox{0.92\textwidth}{
\centering
\texttt{camera / heading -> corridor localizer -> graph planner -> subgoal ->
simple local controller -> optional live loop}
}}
\end{center}

That means the perception/planning/runtime backbone is now real.


\section{What Is Still Missing}

Important things are still missing:

\begin{itemize}[nosep]
  \item final safety integration inside the live loop,
  \item a controller that can follow subgoals without spinning in place,
  \item reliable teleoperation and manual recovery tooling,
  \item cleaner no-path and low-confidence runtime behavior,
  \item final checkpoint sequencing tests on the robot,
  \item and the decision of whether MBRA should replace or augment the simple
        local controller.
\end{itemize}

So the project is \textbf{not finished}.  But it is no longer vague.  It now
has a real technical spine.


\section{What The Next Best Step Is}

The next best step is:

\begin{enumerate}[nosep]
  \item make manual teleoperation reliable,
  \item improve the simple controller so it follows nearby subgoals without
        spinning,
  \item validate short forward targets such as ``current step $\rightarrow$
        current step + 50'',
  \item then decide whether to keep upgrading the heuristic controller or begin
        MBRA integration.
\end{enumerate}

In practical terms:

\begin{quote}
The project has moved from ``can we localize and plan in the corridor?'' to
``how do we make the controller behave correctly on the real robot?''
\end{quote}


\section{Bottom Line}

What we built so far is not random glue code.  It is a deliberate indoor
navigation baseline for a known corridor.

The logic of the system is:

\begin{itemize}[nosep]
  \item build memory of the corridor,
  \item localize against that memory,
  \item plan a route through the graph,
  \item choose a nearby subgoal,
  \item and produce a local motion command.
\end{itemize}

The reason this is the right direction right now is simple:

\begin{itemize}[nosep]
  \item it matches the actual competition setting,
  \item it uses the strongest existing code in the repo,
  \item it is easier to debug than a fully end-to-end approach,
  \item and it keeps the controller layer open for improvement later.
\end{itemize}

So if you want the cleanest summary:

\begin{quote}
We have already built the corridor-specific indoor perception/planning backbone.
The project is now centered on validating the live loop and deciding how much
the controller should remain simple versus how much MBRA should be integrated.
\end{quote}


\section{Related Documents}
\begin{itemize}
  \item \texttt{erc3\_full\_documentation.tex} --- single master guide for the complete project story and current architecture.
  \item \texttt{live\_indoor\_runtime\_story.tex} --- indoor evolution, MBRA integration, and checkpoint-step runtime behavior.
  \item \texttt{live\_outdoor\_ultra\_marathon\_story.tex} --- outdoor and marathon runtime evolution with safety-layer reasoning.
  \item \texttt{outdoor\_perception\_review.tex} --- depth/semantic perception findings and their runtime implications.
\end{itemize}

\end{document}
